The Linear Representation Hypothesis (\lrh) suggests that semantic concepts are encoded as directions in the activation space of Large Language Models (LLMs). While linear manipulations along these directions are effective for related concepts, the behavior of the representation space when interpolating between incongruent or mutually exclusive concepts remains poorly understood. In this work, we investigate the geometry of activation space by performing linear interpolation (\lerp) between incongruent concepts in \gpttwo. We find that such paths exhibit extreme non-Euclidean geometry, characterized by a 480$\times$ spike in the information metric (KL-divergence rate) compared to \congruent paths. Furthermore, interpolation through incongruent regions leads to a 30.7\% increase in Sparse Autoencoder (\sae) reconstruction error and a significant rise in output entropy, indicating that the model's semantic manifold consists of highly curved, disjoint regions. These results suggest a ``curvature crisis'' where linear transitions between unrelated semantic clusters drift off-manifold into unlearned regions of activation space, posing challenges for robust model steering and interpretability.
